 
 
 When a multi-threaded program with a shared memory executes, the order in which memory operations become visible can affect the behavior of the program. The easiest reasoning for such behaviour is to assume that all memory operations occur in a single global order. This assumption is captured by Sequential Consistency (SC)[\cite{SC-Lamport78}]. Under SC, the operations of all processors can be viewed as being interleaved in some sequential order, while the order of operations within each processor is preserved. However, implementing sequential consistency at the hardware level is expensive. Modern processors use several techniques to improve performance, such as store buffers, caches, pipelining, and out-of-order execution. Example, under a memory model called Total Store Order (TSO) \cite{sun1991sparc}, a processor may temporarily keep a store in a buffer while continuing to execute later instructions.  Therfore, the order in which memory operations are observed by different processors may not  always correspond to a single global order. This motivated the development of weak memory models. A weak memory model deviates from SC by setting rules that allow specific memory operations to be observed in an order different from their program order. This encourages hardware optimization while still providing certain guarantees to the programmer.

Several hardware memory models were proposed. Examples include processor consistency (PC) \cite{gharachorloo1990memory}, release consistency (RC) \cite{gharachorloo1990memory}, Weak Ordering \cite{dubois1986memory}, Total Store Order (TSO), Partia Store Order (PSO) \cite{sun1991sparc,sparc1994architecture}.  
Different architecture provided different guarantees. Gharachorloo in 1992  \cite{GHARACHORLOO1992399} summarized the behviours of all these models.

 Besides hardwares, compilers also perform optimizations that can change the order of memory operations thus introducing  additional behaviors that need to be considered. This led to the development of weak memory models extended to programming languages. But weaker model led to more behaviours of the concurrent program. Now it became necessary to check if an observed behaviour is faithful to the permitted model. 
 In 1997, Gibbons and Korach had  studied the problem of testing whether an observed execution of a shared-memory system is consistent with Sequential Consistency. Given the observed reads and writes, their problem asks whether there exists a sequential ordering of these operations that respects the program order of each processor and explains the observed behavior. They showed that this problem is computationally difficult in general, establishing NP-completeness results for several variants. There are several works on testing consistency of SC by varying . \cite{CantinLS05} proved that the problem remains NP-complete even with a single memory
 
 
 
 
location (but with unboundedly many threads), which also implies NP-completness for all memory
models that adhere to the “SC-per-location” property, such as TSO, PSO, RA, SRA and Relaxed 

This led to the development of programming-language memory models, where the language specification itself defines the behavior of concurrent programs. Example include Java Memory Model (JMM), 


. The Java memory model was revised significantly in Java 5, providing a formal description of how threads interact through shared memory. This work highlighted the importance of distinguishing ordinary memory operations from synchronization operations. Similar issues later became central to the design of the C and C++ memory models.


This approach became an important part of the C11 and C++11 memory models. These standards introduced atomic operations with different memory-ordering constraints, including relaxed, acquire, release, acquire-release, and sequentially consistent orderings. Consequently, programmers can select the amount of ordering required for a particular operation instead of having every operation follow sequential consistency.

With the development of formal memory models, the way in which executions are represented has also changed. In earlier discussions, weak memory behavior was often described informally as a processor "reordering" instructions. Modern approaches generally describe an execution using a set of memory events together with relations between those events. For example, program order captures the order of operations within a thread, while a reads-from relation identifies the write from which a read obtains its value. Other relations are used to represent synchronization and ordering constraints. A memory model then specifies which combinations of these events and relations constitute a valid execution.

This formal view is particularly useful for verification. Given an observed execution of a concurrent program, one can ask whether the execution is allowed by a particular memory model. This is the basis of the memory consistency checking problem. Such questions are important because an execution that appears impossible under sequential consistency may nevertheless be valid under a weaker model. Determining whether an observed behavior is consistent with the underlying memory model therefore becomes an important problem in the verification of concurrent programs.

The history of weak memory models can thus be viewed as a gradual movement from a simple global-order view of concurrency towards more flexible and formal descriptions of concurrent executions. Sequential consistency provided a simple starting point, while hardware memory models introduced controlled relaxations to improve performance. Programming-language memory models subsequently had to account for compiler optimizations and synchronization. More recently, formal execution-based models have provided a common framework for studying these behaviors and for developing algorithms to verify them.

Overall, the development of weak memory models reflects a fundamental trade-off between performance and programmability. Stronger memory ordering makes concurrent programs easier to reason about, but it can restrict hardware and compiler optimizations. Weaker ordering provides greater freedom for optimization, but makes reasoning about program behavior more difficult. The study of weak memory models is therefore not only about describing hardware behavior; it is also about finding an appropriate balance between these two requirements.
